\chapter{Background}
\label{ch:background}

\section{The Ramblers Problem}

A terrain map is a grid of heights $H(r,c) \in \{0,\dots,255\}$. A state is a cell $(r,c)$, and the rambler can move to any of its four neighbours. The cost of stepping from cell $u$ to a neighbouring cell $v$ is
\begin{equation}
  c(u,v) =
  \begin{cases}
    1 & \text{if } H(v) \le H(u),\\
    1 + H(v) - H(u) & \text{if } H(v) > H(u),
  \end{cases}
  \label{eq:cost}
\end{equation}
so descending is as cheap as walking on the flat and every metre climbed is paid for. The task is to find a path of least total cost $C^*$ from a start cell to a goal cell.


\section{Search Algorithms}

Both algorithms keep an \emph{open} list of generated nodes and a \emph{closed} list of expanded nodes. A node records its state, its cost so far $g$, and its parent. When a successor reaches a state already on open or closed by a cheaper route, that node is updated and, if it was closed, moved back to open. Search stops when a goal node is \emph{selected} for expansion, not when it is generated, which is what makes the returned path optimal.

\subsection{Branch and bound}
Branch and bound \citep{lawler1966branch} always selects the open node with the lowest $g$. With the duplicate check above it is uniform-cost search, the node-selection form of Dijkstra's algorithm \citep{dijkstra1959note,russell2020aima}.

\subsection{A*}
A* selects the open node with the lowest $f = g + h$, where $h$ estimates the remaining cost to the goal \citep{hart1968formal}. With $h \equiv 0$ it reduces to branch and bound. This was checked in the code: on every instance tested, A* with $h \equiv 0$ reproduced branch and bound's path cost and node count exactly.


\section{Heuristics}

Let $(r,c)$ be the current cell and $(r_g,c_g)$ the goal. Table~\ref{tab:heuristics} lists the five heuristics compared.

\begin{table}[ht]
\centering
\caption{Heuristics compared. ``Admissible'' means the heuristic never exceeds the true remaining cost.}
\label{tab:heuristics}
\begin{tabular}{llc}
\toprule
Name & Definition & Admissible \\
\midrule
Euclidean $\hE$ & $\lfloor \sqrt{(r-r_g)^2 + (c-c_g)^2} \rfloor$ & yes \\
Manhattan $\hM$ & $|r-r_g| + |c-c_g|$ & yes \\
Ascent $\hA$ & $\max\bigl(0,\; H(r_g,c_g) - H(r,c)\bigr)$ & yes \\
Manhattan+Ascent $\hMA$ & $\hM + \hA$ & yes \\
Absolute height $\hAbs$ & $|H(r_g,c_g) - H(r,c)|$ & no \\
\bottomrule
\end{tabular}
\end{table}

\begin{proposition}
$\hMA$ is admissible and consistent, and it dominates $\hM$, $\hE$ and $\hA$.
\end{proposition}
\begin{proof}
Write each step cost in~\eqref{eq:cost} as $1 + \max(0, \Delta)$, where $\Delta$ is the height change of the step. Any path from $u$ to the goal has at least $\hM(u)$ steps. Its positive height changes sum to at least the net rise $H(\text{goal}) - H(u)$, and to at least 0. The path therefore costs at least $\hM(u) + \hA(u) = \hMA(u)$, which proves admissibility. For consistency, take one step from $u$ to $v$. Then $\hM(u) \le 1 + \hM(v)$. The inequality $\max(0, a+b) \le \max(0,a) + \max(0,b)$ gives $\hA(u) \le \max(0, H(v)-H(u)) + \hA(v)$. Adding the two gives $\hMA(u) \le c(u,v) + \hMA(v)$. Dominance holds because $\hA \ge 0$, $\hM \ge 0$ and $\hE \le \hM$.
\end{proof}

$\hAbs$ is the ``height difference'' heuristic of the original version. It also charges for descents, which are cheap. A single downhill step of 100 costs 1, but $\hAbs$ estimates it at 100. By standard results, a consistent heuristic that dominates another expands no more nodes, ignoring ties \citep{pearl1984heuristics,dechter1985generalized}. $\hMA$ is therefore expected to expand the fewest nodes of the admissible heuristics. No such guarantee applies to $\hAbs$.
